% worked_examples.tex — step-by-step runbooks for the checked-in examples/*
% topologies, ordered from the easiest to the most advanced. Companion to,
% not a repeat of, the per-command reference in Text/user_guide.tex and the
% narrative tutorials under docs/tutorials/: this section gives the exact
% command sequence for each topology and explains what changes between
% them, and points to those other documents for flag semantics and full
% expected-output transcripts.
% -----------------------------------------------------------------------
\section{Worked Examples: From \texttt{CGSil1} to \texttt{cawaqsviz}}
\label{sec:worked-examples}

\cgspkg{} ships four checked-in \texttt{examples/} topologies, and a
same-numbered tutorial under \texttt{docs/tutorials/} for each level below
(\texttt{docs/tutorials/README.md} lists all three in order). Run through
them in order: each one adds exactly one new idea on top of the last, so
together they form a guided path from the simplest possible \texttt{.cgs}
to a real project with no \texttt{.cgs} of its own at all.

\textbf{Every command below is a Pixi task.} \cgscli{} is not a globally
installed executable --- always run \texttt{pixi run cgitsync ...}, never a
bare \texttt{cgitsync ...} (see the \S\ref{chap:user-guide} ``CLI
Overview'' note above).

\begin{center}
\begin{tabular}{lllp{5.3cm}}
\textbf{Level} & \textbf{Topology} & \textbf{Entry point} & \textbf{New idea introduced} \\
\hline
1 & \texttt{CGSil1} & hand-authored \texttt{.cgs} & minimal spec, name-matching auto-mount \\
2 & \texttt{cawaqs} & hand-authored \texttt{.cgs} & the same minimal-field style as Level~1, at 19-repo scale \\
3a & \texttt{cawaqsviz} & hand-authored \texttt{.cgs} & explicit overrides on a real, irregular tree \\
3b & \texttt{cawaqsviz} & \texttt{discover} & drafting a \texttt{.cgs} from a checkout, no authoring at all \\
3c & \texttt{cawaqsviz} & \texttt{import-submodules} & migrating an existing git-submodules project \\
\end{tabular}
\end{center}

Levels~3a--3c are not three different projects: they are three different
\emph{starting points} for the same real project, \texttt{cawaqsviz}, chosen
so this section can demonstrate every route \cgspkg{} offers for arriving
at a working \texttt{.cgs} --- write it by hand, derive it from a checkout,
or migrate it from an existing submodules setup --- without inventing a
fourth toy topology to do it. It is the most advanced level not because
\texttt{cawaqsviz} is a larger tree than \texttt{cawaqs} (it is much
smaller, three repos against nineteen) but because, unlike either earlier
level, it has no \texttt{.cgs} anywhere upstream to copy from.

% -----------------------------------------------------------------------
\subsection{Level 1 --- \texttt{CGSil1}: the minimal synthetic topology}

\texttt{CGSil1} is the easiest way into \cgspkg{}: a 4-repository,
mixed-provider tree with every field left at its default. Its full
\texttt{.cgs} and the rationale for each line are in \S\ref{chap:user-guide}
``Local Authoring Spec'' above;
\texttt{docs/tutorials/01\_first\_multi\_repo\_workspace.md} carries the
full transcript with expected output for every step. What follows is the
command sequence itself.

\begin{lstlisting}[language=bash]
# 1. Fetch the root repo and cgitsync itself (nested mode: cgitsync is
#    cloned inside the tree it will manage).
$ git clone https://gitlab.com/CGS_test/CGSil1
$ cd CGSil1 && git clone https://github.com/flipoyo/ComplexGitSync
$ cd ComplexGitSync

# 2. Check the spec offline, no network beyond parsing.
$ pixi run cgitsync validate ../CGSil1.cgs

# 3. Preview the topology before touching disk.
$ pixi run cgitsync view-tree ../CGSil1.cgs

# 4. Clone every declared repo and reach READY.
$ pixi run cgitsync initialise ../CGSil1.cgs

# 5. The ordinary git cycle, run tree-wide from here on --- no path
#    argument needed, the .gts snapshot is discovered automatically.
$ pixi run cgitsync pull
$ pixi run cgitsync add
$ pixi run cgitsync commit "my commit message"
$ pixi run cgitsync push

# 6. Minimalist release: stage, commit, pull, push, freeze in one call.
$ pixi run cgitsync freeze-release v1.1.0 "release v1.1.0"

# 7. Return to a frozen release at any later point.
$ pixi run cgitsync launch-release v1.1.0
\end{lstlisting}

The only \texttt{.cgs} idea this level relies on is the
\textbf{name-matching auto-mount}: because \texttt{CGSil1}'s repository
identifier's last segment equals \texttt{project = "CGSil1"}, the root is
inferred at \texttt{relative\_path = "."} with no override written anywhere.
That convenience is exactly what stops being safe once a project's
identifier and display name diverge --- which is where Level~3a picks up.

% -----------------------------------------------------------------------
\subsection{Level 2 --- \texttt{cawaqs}: the same minimal style, at scale}

\texttt{cawaqs} scales Level~1's minimal, no-override authoring style up to
a real 19-repository build tree (root + 17 physics libraries across five
GitLab groups + one required build-tooling repo), instead of a synthetic
4-repo one. Every entry's on-disk layout already matches the default
\texttt{relative\_path = <repo\_name>}, so not a single \texttt{relative\_path}
override is written anywhere in \texttt{examples/cawaqs.cgs} --- only the
shared branch-fallback convention. Full narrative in
\texttt{docs/tutorials/02\_onboarding\_a\_real\_build\_tree.md}.

\begin{lstlisting}[language=bash]
$ pixi run cgitsync validate examples/cawaqs.cgs
$ pixi run cgitsync bootstrap examples/cawaqs.cgs cawaqs-smoke-test
$ pixi run cgitsync view-tree --search-dir <path bootstrap printed>

# Move the whole 19-repo tree to a shared branch in one call, falling
# back to main per-repo wherever that branch does not exist.
$ pixi run cgitsync checkout my-feature-branch
\end{lstlisting}

Reaching \texttt{READY} here only replaces the \emph{repository-fetching and
branch-selection} half of \texttt{cawaqs}'s own
\texttt{make\_Cawaqs.sh}/\texttt{make\_Cawaqs\_from\_branches.sh} scripts.
\cgspkg{} does not compile anything: the remaining
\texttt{make -f Makefile} build step is still the user's job, and only
succeeds from this tree with \texttt{LIB\_HYDROSYSTEM\_PATH} left unset or
pointed at the bootstrapped root --- the alternative shared, decoupled
install layout those scripts also support is out of scope for \cgspkg{}'s
containment model.

\texttt{discover} (Level~3b below) is not a route into \texttt{cawaqs.cgs}:
the 17 library repositories never coexist inside one directory tree until
\emph{after} a \texttt{.cgs} already lists them, so there is no single
checkout for a filesystem walk to scan. Authoring this one from documented
developer knowledge, the way Level~3a does for \texttt{cawaqsviz}, is the
only available route.

% -----------------------------------------------------------------------
\subsection{Level 3 --- \texttt{cawaqsviz}: three ways to the same \texttt{.cgs}}

\texttt{cawaqsviz} (\url{https://gitlab.com/cawaqs/gviz/cawaqsviz}) is a
real GitLab project, nested three path segments deep, with two GitHub
children that were, until recently, plain git submodules. It has no
\texttt{.cgs} of its own anywhere upstream. That makes it a good vehicle for
showing all three ways \cgspkg{} can produce a working \texttt{.cgs} for a
project it does not already control the source of, ordered here by how much
the operator has to already know about the tree before running the command.
Full narrative, as one combined ten-step procedure rather than three
independent routes, in \texttt{docs/tutorials/03\_adopting\_a\_real\_project.md}.

\subsubsection{3a --- Write it by hand}

The most explicit, and most error-prone, route: read the project's own
\texttt{.gitmodules} and topology, then author \texttt{examples/cawaqsviz.cgs}
directly. Its corrected content and the exact mistakes the first draft made
(wrong nested-subgroup identifier, missing \texttt{relative\_path}
overrides) are already spelled out in \S\ref{chap:user-guide} ``Local
Authoring Spec''; this level is about running it, not re-deriving it:

\begin{lstlisting}[language=bash]
$ pixi run cgitsync validate examples/cawaqsviz.cgs
$ pixi run cgitsync bootstrap examples/cawaqsviz.cgs cawaqsviz-demo
$ pixi run cgitsync view-tree --search-dir <path bootstrap printed>
\end{lstlisting}

\texttt{bootstrap} is used instead of \texttt{initialise} here because, unlike
\texttt{CGSil1} above, this example is not meant to have \cgspkg{} cloned
inside the tree it manages --- see \texttt{bootstrap} in
\S\ref{chap:user-guide} for why the two commands pick different default
locations.

\subsubsection{3b --- Derive it with \texttt{discover} (autodiscover mode)}

The opposite route: start from nothing but a checkout, and let \cgspkg{}
read the filesystem instead of writing the \texttt{.cgs} by hand. This is
the route to prefer whenever a checkout is already available, since it
cannot repeat the hand-authoring mistakes of 3a --- \texttt{relative\_path}
falls out of the walk directly instead of being typed.

\begin{lstlisting}[language=bash]
# A plain clone leaves submodule paths empty; initialise them first so
# discover has something to find at those paths. --init only (not
# --recursive): a nested submodule one level deeper would otherwise show
# up as a fourth repository discover has no use for here. (Want that
# nested submodule too? Use --init --recursive here and
# import-submodules --recursive --apply in 3c.)
$ git clone https://gitlab.com/cawaqs/gviz/cawaqsviz cawaqsviz
$ cd cawaqsviz && git submodule update --init

# Dry run: report what discover sees, write nothing.
$ pixi run cgitsync discover . --max-depth 3

# Satisfied with the report: draft the .cgs.
$ pixi run cgitsync discover . --write cawaqsviz-discovered.cgs
$ pixi run cgitsync validate cawaqsviz-discovered.cgs
\end{lstlisting}

The draft reconstructs 3a's hand-corrected file field-for-field: the root at
\texttt{relative\_path = "."} and both children at their real submodule
paths (\texttt{external/HydrologicalTwinAlphaSeries},
\texttt{docs/CWV\_user\_guide}) rather than their bare repo names. Neither
child has a \texttt{.cgs} of its own, but \texttt{discover} writes no
\texttt{nested\_config} override for that: the default \texttt{"auto"}
already resolves a repository with zero nested \texttt{*.cgs} matches as a
normal leaf. See \texttt{discover} in \S\ref{chap:user-guide} for exactly
what is and is not reported (no network calls, unresolvable remotes are
warned about and left out, never guessed at).

\subsubsection{3c --- Migrate it with \texttt{import-submodules}}

The historical route: before either \texttt{.cgs} above existed,
\texttt{cawaqsviz} tracked both children as real git submodules. This mode
starts from that state and converts it, rather than describing a tree that
is assumed to already be plain clones.

\begin{lstlisting}[language=bash]
# Dry run against a real cawaqsviz checkout that still has submodules:
# reports name, path, URL, branch per submodule; changes nothing.
$ pixi run cgitsync import-submodules /path/to/cawaqsviz

# Convert: drop the gitlinks, retire .gitmodules, extend .gitignore.
$ pixi run cgitsync import-submodules /path/to/cawaqsviz --apply

# Review and commit the resulting working-tree changes as usual.
$ cd /path/to/cawaqsviz && git status
$ git commit -m "chore: retire git submodules in favour of ComplexGitSync"
\end{lstlisting}

This is the whole job: turning gitlinks into plain clones on disk. It does
not also write a \texttt{.cgs} --- \texttt{.gitmodules} never records the
root's own identity, and a checkout worth converting already has a
\texttt{.cgs} (3a) or can get one from \texttt{discover} (3b), run before
or after \texttt{--apply}. The full before/after picture (index,
\texttt{.gitmodules}, \texttt{.gitignore}) and the automated test that
proves it against fixture repositories are in
\texttt{docs/tutorials/03\_adopting\_a\_real\_project.md} \S4 --- this
level is the condensed command sequence, not a restatement of that
walkthrough. \texttt{import-submodules} and \texttt{discover} answer
different questions about the same tree: \texttt{discover} reports what is
\emph{checked out}, \texttt{import-submodules} reads what git
\emph{declares} in \texttt{.gitmodules} and acts on it --- a submodule path
that was never initialised is invisible to 3b but not to 3c.

\subsubsection{A shared caveat across all three}

\cgspkg{} stores no credentials and has no private-repository
authentication story (\S\ref{chap:user-guide} ``\texttt{import-submodules}''
and ``\texttt{discover}''). All three routes above assume every repository
in the tree is reachable anonymously or through credentials the ambient
environment (\texttt{ssh-agent}, an HTTPS credential helper) already holds.
